% TOP_PROBLEM: {"id": 5200015, "title": "Open Problems on Billiards and Geometric Optics", "category_id": 11, "category": {"id": 11, "name": "dynamical_systems", "display_name": "Dynamical Systems", "description": "Problems about long-term behavior of deterministic systems, Hamiltonian dynamics, and periodic orbits.", "slug": "dynamical-systems", "order_index": 11, "created_at": "2026-07-31T15:26:25.671Z"}, "status": "partially_solved"}
% FIRST_POSTED: null
% ENTRY_KIND: statement_only
% Imported from: batch03.tex; UnsolvedMath ID 5200015
\documentclass[11pt]{article}
\usepackage[margin=1in]{geometry}
\usepackage{amsmath,amssymb,mathtools}
\usepackage{fontspec,unicode-math}
\setmainfont{Latin Modern Roman}
\setsansfont{Latin Modern Sans}
\setmathfont{Latin Modern Math}
\usepackage{xurl,hyperref}
\hypersetup{colorlinks=true,urlcolor=blue,linkcolor=blue}
\newcommand{\sref}[2]{\hyperlink{source:#1:#2}{[S#2]}}
\newcommand{\cataloguescope}[1]{\paragraph{Recorded mathematical question.}#1}
\begin{document}
% UnsolvedMath ID 5200015; source catalogue code AMR-051-0015
\setcounter{section}{5200014}
\section{Open Problems on Billiards and Geometric Optics}\label{5200015}
{\small\sffamily UnsolvedMath ID 5200015 · Dynamical Systems}
\subsection{Definitions and mathematical statement}

\paragraph{Shared notation.}$\mathbb N=\{1,2,\ldots\}$, and $\mathbb Z,\mathbb Q,\mathbb R,\mathbb C$ denote the integers, rationals, reals and complex numbers. Any other conventions are specified locally.


An oval is a smooth closed strictly convex plane curve. The phase points are oriented chords $(x,y)$ and the billiard map sends $(x,y)$ to $(y,z)$ according to the parallel-tangent rule. For a polygon, use the line containing the edge at a regular impact point; trajectories hitting a vertex or an undefined reflection are excluded from regular orbits. A periodic orbit returns to its initial chord after finitely many steps. A stadium consists of two equal semicircular caps joined by parallel straight segments. The source asks broadly about chaos without selecting a particular numerical criterion.
\paragraph{Statement recorded in the supplied source.}
In a planar symplectic billiard on an oval, the chord $xy$ reflects to $yz$ when the tangent at $y$ is parallel to $xz$; define polygonal symplectic billiards similarly. (1) Classify polygons for which every symplectic billiard orbit is periodic. (2) Does every polygon have a periodic orbit? (3) Is the symplectic billiard in a stadium chaotic? \sref{5200015}{2}
\subsection{Short English statement}
Classify polygons with all regular symplectic billiard orbits periodic, decide whether every polygon has some periodic orbit, and establish whether the symplectic stadium dynamics is chaotic.
\subsection{Sources}
\begin{description}
\item[\hypertarget{source:5200015:1}{S1}] ULAM.AI and the UnsolvedMath Contributors, \emph{UnsolvedMath: A Curated Collection of Open Mathematics Problems}, record 5200015 (AMR-051-0015). CC BY 4.0. \url{https://huggingface.co/datasets/ulamai/UnsolvedMath}.
\item[\hypertarget{source:5200015:2}{S2}] M. Bialy, C. Fierobe, A. Glutsyuk, M. Levi, A. Plakhov and S. Tabachnikov, \emph{Open problems on billiards and geometric optics} (2021), arXiv:2110.10750v2, Serge Tabachnikov, Problem 3, pp. 11--12. \url{https://arxiv.org/pdf/2110.10750v2}.
\end{description}
\label{end:5200015}
\end{document}
